\section{扰动作用量泛函及其作用谱}
\label{sec:perturbed-action-functionals-spectrum}

给定 1-周期哈密顿量 $H:(\mathbb R/\mathbb Z)\times M\to\mathbb R$ 后，考虑扰动泛函
$\mathcal A_H:\widetilde{\mathcal L}(M)\to\mathbb R$：
\begin{equation}
 \mathcal A_H([\gamma,w])
 =\mathcal A_0([\gamma,w])-\int H(t,\gamma(t))\,dt
 =-\int w^*\omega-\int H(t,\gamma(t))\,dt.
 \label{eq:18-perturbed-action-functional}
\end{equation}
除非另作说明，我们总取 1-周期\emph{均值归一化哈密顿量}，即满足
\begin{equation}
 \int_M H_t\,\omega^n=0\qquad\text{对所有 }t\in S^1.
 \label{eq:18-mean-normalized-hamiltonian}
\end{equation}
记 $\mathcal H_0(M)=C_0^\infty(S^1\times M)$ 为这类均值归一化哈密顿量的集合。

\begin{lemma}
\label{lem:18-critical-points-perturbed-action}
$\mathcal A_H$ 的临界点集为
\[
 \operatorname{Crit}(\mathcal A_H)
 =\{[z,w]\mid z\in\operatorname{Per}(H),\ \partial w=z\},
\]
而 $\Gamma$ 在 $\widetilde{\mathcal L}_0(M)$ 上的作用典范地限制到该集合上。
\end{lemma}

\begin{definition}
\label{def:18-action-spectrum}
\stfhIndex{I027}{Zuoyongpu}{作用谱（action spectrum）}
定义 $H$ 的\emph{\termActionSpectrum{}}为
\[
 \operatorname{Spec}(H):=\{\mathcal A_H([z,w])\in\mathbb R\mid
 [z,w]\in\widetilde{\mathcal L}_0(M),\ z\in\operatorname{Per}(H)\},
\]
即 $\mathcal A_H:\widetilde{\mathcal L}(M)\to\mathbb R$ 的临界值集合。对每个给定的
$z\in\operatorname{Per}(H)$，记
\[
 \operatorname{Spec}(H;z)=\{\mathcal A_H([z,w])\in\mathbb R\mid
 [z,w]\in\pi^{-1}(z)\}.
\]
\end{definition}

\stfhIndex{I009}{00-gammaomega-torsor}{$\Gamma_\omega$-挠子（$\Gamma_\omega$-torsor）}
注意，$\operatorname{Spec}(H;z)$ 是一个 $\Gamma_\omega$-\termTorsor{}
\footnote{译者注：原书称固定 $z$ 的 $\operatorname{Spec}(H;z)$ 是“位于
$\operatorname{Per}(H)$ 上的 $\Gamma_\omega$ 主丛”，但固定纤维不能同时是整个基底
上的丛；译文按随后不交并公式表述为位于 $z$ 上的挠子。}，位于
$z\in\operatorname{Per}(H)$ 之上。因此
\[
 \operatorname{Spec}(H)=\bigsqcup_{z\in\operatorname{Per}(H)}
 \operatorname{Spec}(H;z).
\]
由于 $\Gamma_\omega=\omega(\pi_2(M))$ 是 $\mathbb R$ 的可数子群，它要么离散，要么
在 $\mathbb R$ 中稠密。下面的结论是研究改变哈密顿量 $H$ 时
$\mathcal A_H$ 的临界值行为的重要工具。这里的证明取自 (Oh02)。

\begin{lemma}
\label{lem:18-action-spectrum-measure-zero}
设 $H$ 是任意周期哈密顿量。则对任意 $H$，$\operatorname{Spec}(H)$ 都是
$\mathbb R$ 的零测集。
\end{lemma}

\begin{proof}
首先，对每个给定的周期轨道 $z$，$\operatorname{Spec}(H;z)\subset\mathbb R$ 是可数
集，因为 $\Gamma_\omega$ 可数。在每个 $z\in\operatorname{Per}(H)$ 的一个管状邻域
中考虑\personPoincare{}回归映射。更确切地说，取 $z(0)$ 的一个小邻域 $V\subset M$，将
$V$ 与以 $z(0)$ 为球心的 $2n$ 维球 $B^{2n}(\delta)$ 认同。再取球邻域
$V'=B^{2n}(\delta')$，满足 $V\Subset V'$，使（第一）\personPoincare{}回归映射
\[
 R_z:V\longrightarrow V',\qquad p\longmapsto\phi_H^1(p)
\]
有定义。现在如下定义从 $V$ 到 $S^1\cong\mathbb R/\mathbb Z$ 到 $M$ 的分片光滑映射
空间的连续映射：对每个 $p\in V$，先沿 $X_H$ 的流前进，再在 $V'$ 与
$B^{2n}(\delta')$ 的认同下沿直线从 $R_z(p)$ 返回 $p$。最后缩放环路的定义域，将其
重新参数化为 $[0,1]$。

记 $z_p$ 为上述构造中与 $p\in V$ 对应的环路，$\mathcal V_z\subset\mathcal L_0(M)$
为赋值 $p\mapsto z_p$ 的像。显然 $z_p$ 与 $z$ 同伦，因此任何以 $z$ 为边界的给定
圆盘 $w$ 都可自然延拓成以 $z_p$ 为边界的圆盘，记为 $w_p$。容易验证函数
\[
 h:\pi^{-1}(\mathcal V_z)\longrightarrow\mathbb R,
 \qquad h([z_p,w_p]):=\mathcal A_H([z_p,w_p])
\]
在 $\pi^{-1}(\mathcal V_z)$ 上光滑，并且它的临界值包含
$\operatorname{Spec}(H;z)$ 附近的 $\mathcal A_H$ 的临界值。这可通过写出
$\mathcal A_H([z_p,w_p])$ 的显式表达式并作简单的局部计算来证明。注意
$\pi^{-1}(\mathcal V_z)$ 是有限维（事实上为 $2n$ 维）流形，由\personSard{}定理，临界值集
在 $\mathbb R$ 中是零测集。所有 $z\in\operatorname{Per}(H)$ 的初始点 $z(0)$ 组成
紧集，因此有限多个这样的管状邻域连同其补集覆盖 $M$。故
$\operatorname{Spec}(H)\subset\mathbb R$ 包含在有限个零测集的并中，从而本身也是
零测集。
\end{proof}

\begin{exercise}
\label{ex:18-smoothness-local-action-function}
证明上面证明中定义的函数 $h$ 确实光滑。
\end{exercise}

注意，当 $H=0$ 时，
\[
 \operatorname{Spec}(H)=\Gamma_\omega.
\]

\begin{definition}
\label{def:18-homotopic-hamiltonians}
若 $\phi_H^1=\phi_F^1$，并且相伴的哈密顿路径
$\phi_H,\phi_F\in\mathcal P(\operatorname{Ham}(M,\omega),\mathrm{id})$ 保端点道路同伦，
则称两个哈密顿量 $H$ 与 $F$ \emph{同伦}。此时记 $H\sim F$，并以
$\widetilde{\operatorname{Ham}}(M,\omega)$ 表示等价类的集合。
\footnote{译者注：原书在定义 18.3.5 中把第二条相伴路径印作 $\phi_K$，但定义的两个
哈密顿量是 $H$ 与 $F$；译文按上下文改为 $\phi_F$。}
\end{definition}

下面的命题是研究哈密顿量形变下作用量泛函临界值的另一重要工具。它在非球面情形
由 (Schw00)、(Po01) 证明，在一般情形由 (Oh05a) 证明。

\begin{proposition}
\label{prop:18-homotopic-hamiltonians-same-spectrum}
设 $F,G$ 均为均值归一化哈密顿量。若 $F\sim G$，则作为 $\mathbb R$ 的子集，
\[
 \operatorname{Spec}F=\operatorname{Spec}G.
\]
\end{proposition}

证明前先作些准备。设 $\phi$ 是给定的哈密顿微分同胚，
$\{F^s\}_{s\in[0,1]}$ 是一条哈密顿量道路，满足
$F^0=G$、$F^1=F$，且对所有 $s\in[0,1]$ 都有 $\phi_{F^s}^1=\phi$。考虑哈密顿
微分同胚的二参数族
\[
 \phi(s,t)=\phi_{F^s}^t\circ(\phi_G^t)^{-1}.
\]
记 $F(s,t,x)=F^s(t,x)$，并令 $K=K(s,t,x)$ 是 $s$ 方向上的哈密顿函数，即与向量场
\[
 \frac{\partial\phi_{F^s}^t}{\partial s}\circ(\phi_{F^s}^t)^{-1}
\]
相伴的哈密顿量。若要求 $K(s,t,\cdot)$ 归一化，则 $K$ 唯一，而这种归一化总能实现。
此外，由 (2.4.26)，
\[
 \frac{\partial X_F}{\partial s}-\frac{\partial X_K}{\partial t}
 +[X_K,X_F]=0,
\]
继而由 (2.2.9) 和引理 2.3.4，
\[
 \frac{\partial F}{\partial s}-\frac{\partial K}{\partial t}-\{K,F\}=c(s,t),
\]
其中 $c$ 只依赖于 $(s,t)$。由于 $F$ 与 $K$ 都已归一化，而由\personLiouville{}引理，\personPoisson{}括号 $\{K,F\}$ 自动归一化，所以 $c\equiv0$。换言之，$F,K$ 满足
\begin{equation}
 \frac{\partial F}{\partial s}-\frac{\partial K}{\partial t}-\{K,F\}=0.
 \label{eq:18-hamiltonian-two-parameter-identity}
\end{equation}
值得把下面的等价恒等式单列出来。

\begin{lemma}
\label{lem:18-equivalent-two-parameter-identity}
等式~\eqref{eq:18-hamiltonian-two-parameter-identity} 等价于
\begin{equation}
 \frac{\partial F}{\partial s}(s,t,\phi_{F^s}^t(p))
 =\frac\partial{\partial t}\bigl(K(s,t,\phi_{F^s}^t(p))\bigr).
 \label{eq:18-hamiltonian-total-derivative}
\end{equation}
\end{lemma}

\begin{proof}
假设~\eqref{eq:18-hamiltonian-two-parameter-identity} 成立，则直接计算
\begin{align*}
 \frac{\partial F}{\partial s}(s,t,\phi_{F^s}^t(p))
 &=\frac{\partial K}{\partial t}(s,t,\phi_{F^s}^t(p))
   +\{K,F\}(s,t,\phi_{F^s}^t(p))\\
 &=\frac{\partial K}{\partial t}(s,t,\phi_{F^s}^t(p))
   +dK(X_F)(\phi_{F^s}^t(p))\\
 &=\frac\partial{\partial t}\bigl(K(s,t,\phi_{F^s}^t(p))\bigr),
\end{align*}
这就证明了~\eqref{eq:18-hamiltonian-total-derivative}。反过来，只需从右端的最后一行
倒着读回第一行即可。证明完毕。
\end{proof}

现在可以证明命题~\ref{prop:18-homotopic-hamiltonians-same-spectrum}。

\begin{proof}[命题~\ref{prop:18-homotopic-hamiltonians-same-spectrum} 的证明]
设 $\phi$ 是给定的哈密顿微分同胚，取满足上述条件的道路
$\{F^s\}_{s\in[0,1]}$，并使用二参数族
$\phi(s,t)=\phi_{F^s}^t\circ(\phi_G^t)^{-1}$ 及哈密顿量 $F,K$。注意，对每个
$s\in[0,1]$，$t\mapsto\phi(s,t)$ 都定义一条满足 $\phi(s,0)\equiv\mathrm{id}$ 的
可缩哈密顿环路，这是关于 $\{F^s\}$ 的假设所致。记此环路为 $h^s$，记其满足
$\widetilde h^s(0)=\mathrm{id}$ 的提升为 $\widetilde h^s$，它位于
$\widetilde{\operatorname{Ham}}(M,\omega)$ 中。环路 $\widetilde h^s$ 给出一一对应
\footnote{译者注：原书把下式的目标写作 $\operatorname{Crit}\mathcal A_F$；但
$h^s$ 把 $G$ 的轨道送到 $F^s$ 的轨道，且紧接着定义的 $\chi(s)$ 也使用
$\mathcal A_{F^s}$，故译文补上上标 $s$。}
\[
 (\widetilde h^s)_*: \operatorname{Crit}\mathcal A_G
 \longrightarrow\operatorname{Crit}\mathcal A_{F^s},
\]
其公式为
\[
 \widetilde h^s\cdot[z,w]=[h^s\cdot z,\widetilde h^s\cdot w].
\]
因此只需证明函数
\[
 \chi:[0,1]\to\mathbb R,\qquad
 \chi(s)=\mathcal A_{F^s}(\widetilde h^s\cdot[z,w])
\]
为常数。

由于对所有 $s\in[0,1]$ 都有
$\phi_{F^s}^1\equiv\phi_G^1$、$\phi_{F^s}^0\equiv\mathrm{id}$，故对每个 $s$，
$K(s,1,x)$ 和 $K(s,0,x)$ 都是常数；由归一化条件进一步得到
\begin{equation}
 K(s,1,\cdot)=K(s,0,\cdot)\equiv0.
 \label{eq:18-k-endpoint-normalization}
\end{equation}
对 $\chi(s)$ 求导：
\[
 \chi'(s)=d\mathcal A_{F^s}(\widetilde h^s\cdot[z,w])
 \left(\frac\partial{\partial s}(\widetilde h^s\cdot[z,w])\right)
 -\int_0^1\frac{\partial F}{\partial s}
 (s,t,(h^s\cdot z)(t))\,dt.
\]
因为 $\widetilde h^s\cdot[z,w]$ 是 $\mathcal A_{F^s}$ 的临界点，第一项为零。对第二项，
注意 $z$ 是 $\dot x=X_G(t,x)$ 的周期解，故存在 $\phi_G^1$ 的不动点 $p\in M$，使
$z(t)=\phi_G^t(p)$。于是
\[
 (h^s\cdot z)(t)=h_t^s(z(t))
 =(\phi_{F^s}^t\circ(\phi_G^t)^{-1})\circ\phi_G^t(p)
 =\phi_{F^s}^t(p).
\]
从而
\[
 \int_0^1\frac{\partial F}{\partial s}(s,t,(h^s\cdot z)(t))\,dt
 =\int_0^1\frac{\partial F}{\partial s}(s,t,\phi_{F^s}^t(p))\,dt.
\]
代入~\eqref{eq:18-hamiltonian-total-derivative}，在 $[0,1]$ 上积分全导数，并使用
\eqref{eq:18-k-endpoint-normalization}，得到
\[
 \int_0^1\frac{\partial F}{\partial s}(s,t,(h^s\cdot z)(t))\,dt
 =K(s,1,\phi_{F^s}^1(p))-K(s,0,\phi_{F^s}^0(p))=0.
\]
因此 $\chi$ 为常数，命题得证。
\footnote{译者注：原书本式末端两处写成未定义的 $f_1^s(p)$、$f_0^s(p)$；译文按本证明
始终使用的哈密顿流补为 $\phi_{F^s}^1(p)$、$\phi_{F^s}^0(p)$。}
\end{proof}

这个命题说明，归一化哈密顿量 $H$ 的作用谱 $\operatorname{Spec}H$ 只依赖于它的
同伦类。\footnote{译者注：原书此处称“这个引理”，但给出谱不变性结论的是命题
18.3.6；引理 18.3.7 只是其证明中使用的二参数恒等式，译文据此改为“这个命题”。}

\begin{definition}
\label{def:18-spectrum-hamiltonian-path-class}
设 $h\in\widetilde{\operatorname{Ham}}(M,\omega)$。取满足
$h=[\phi,F]$ 的某个（因而任意一个）归一化哈密顿量 $F$，定义 $h$ 的谱为
\[
 \operatorname{Spec}(h):=\operatorname{Spec}F.
\]
\end{definition}

显然，无论 $H$ 是否归一化，作用量泛函 $\mathcal A_H$ 都有定义，其作用谱亦然。记
相应的归一化哈密顿量为 $\underline H$，即
\begin{equation}
 \underline H(t,x)=H(t,x)-\frac1{\operatorname{vol}_\omega(M)}
 \int_M H_t(x)\,d\mu_\omega,
 \label{eq:18-hamiltonian-normalization}
\end{equation}
则\footnote{译者注：原书式 (18.3.21) 的右端写成
$\operatorname{Spec}(H)-\operatorname{Cal}(H)$。但由式 (18.3.15) 的约定
$\mathcal A_H=\mathcal A_0-\int H\,dt$ 及式 (18.3.20) 的
$\underline H=H-\operatorname{Cal}$，直接代入得到下式中的加号；译文据此改正。}
\begin{equation}
 \operatorname{Spec}(\underline H)
 =\operatorname{Spec}(H)+\operatorname{Cal}(H),
 \label{eq:18-spectrum-normalization-shift}
\end{equation}
其中
\[
 \operatorname{Cal}(H):=\frac1{\operatorname{vol}_\omega(M)}
 \int_0^1\int_M H_t(x)\,d\mu_\omega\,dt.
\]
